% !PRINTED 221
\weberpage{221}

The congruence (9) possesses, if $k$ is the number of the mutually distinct primes dividing $n$,
$2^k$ incongruent solutions, and this is therefore the order of the group
$\mathfrak{B}_0$\footnote{Cf. \gesperrt{Dirichlet-Dedekind}, Vorlesungen über Zahlentheorie, 4.\ Aufl., § 37.}.

A rational function $\xi$ of the roots of \gesperrt{one row}, say of the row $R_{1,0}$, can by (1) be
represented rationally by \gesperrt{one of these} roots. If this function has the property of remaining
unchanged when this one root is replaced by another of the same row, that is, if for example
$\xi$ is a \gesperrt{symmetric function of the roots of one row}, then $\xi$ takes, through application of
the substitutions of $\mathfrak{A}$ (or also of $\mathfrak{B}$), only
\weq{10}{\nu = \psi(n)}
different values
\weq{11}{\xi_1, \xi_2, \ldots, \xi_\nu,}
% Source reading \xi_\nu retained: the damaged subscript is fixed by the immediately following count of \nu quantities.
\gesperrt{and if these values are all different from one another, then $\xi$ belongs to the group
$\mathfrak{A}_{1,0}$, and all other symmetric functions of the roots of $R_{1,0}$ can be
represented rationally by $\xi$} (Bd.\ I, § 162).

The $\nu$ quantities (11) are the roots of an irreducible rational equation of degree $\nu$,
which we call a \gesperrt{transformation equation}.

The transformation equation itself has rational numerical coefficients, provided that only
rational numerical coefficients occur in $\xi$. \gesperrt{But it remains irreducible even when $n$th roots
of unity, or indeed any constants whatever, are adjoined}, as follows from the fact that by
substitutions of the group $\mathfrak{B}$ every root of the division equation can be carried
over into every other, and hence also every row into every other row.

By the adjunction of a root $\xi$ of the transformation equation, say of the one belonging to
the group $\mathfrak{A}_{1,0}$, the group of the division equation reduces to
$\mathfrak{A}_{1,0}$; but this latter group is no longer transitive, and instead permutes among
themselves the $\varphi(n)$ roots $x_{h,0}$, in which $h$ runs through a complete system of
incongruent numbers coprime to $n$. The division equation therefore becomes reducible and has a
factor of degree $\varphi(n)$,
%CONT>>
